\title{Parameter-Efficient Orthogonal Finetuning via Factorization}

\begin{abstract}
With the proliferation of large foundation models, retraining such architectures from the ground up has become financially and computationally prohibitive. Consequently, developing methods that allow for efficient adaptation of these models to specialized downstream tasks has gained significant relevance. This work focuses on Orthogonal Finetuning (OFT), a theoretically motivated framework for adapting pre-trained models. While OFT exhibits strong generalization across tasks, its reliance on high-dimensional orthogonal matrices results in a considerable count of trainable parameters. To improve parameter-efficiency, we re-examine OFT through the lens of information transmission and highlight several important requirements for more effective adaptation. Taking inspiration from the efficiency of the Cooley-Tukey fast Fourier transform algorithm, we introduce a new orthogonal parameterization that leverages butterfly structures. Incorporating this construction into OFT, we propose Orthogonal Butterfly~(BOFT), a novel finetuning strategy that substantially reduces parameter usage. BOFT generalizes OFT, which it encompasses as a particular instance, thus providing a broader framework for orthogonal finetuning. Our empirical analysis demonstrates the efficacy of this method by adapting state-of-the-art vision transformers, large language models, and text-to-image diffusion models for a wide array of vision and language downstream tasks.
\end{abstract}

\section{Introduction}

Models such as ChatGPT~\cite{brown2020language,chen2021evaluating} and Stable Diffusion~\cite{rombach2022high} exemplify the outstanding ability of large foundation models to generalize across a diverse range of applications. The advances in performance seen in these models have been accompanied by an unprecedented expansion in their parameter sizes

(\emph{e.g.}, GPT-3 is comprised of roughly 175B parameters). Training such models from scratch is now a daunting challenge for most practitioners, necessitating alternative strategies for advancing the field. Progress increasingly hinges on the ability to effectively repurpose or adapt these existing models for new, task-specific applications, without restarting the training process. There are three principal approaches to such efficient task adaptation: \emph{model finetuning}~\cite{hulora2022,zaken2022bitfit,guo2020parameter,raffel2020exploring,qiu2023controlling,zhang2022adaptive,chavan2023one}, which updates a select portion of the model's parameters while keeping its architecture intact; \emph{adapter tuning}~\cite{rebuffi2017learning,houlsby2019parameter,pfeiffer2020adapterfusion,lin2020exploring,he2021towards}, where new trainable modules are introduced in addition to the base model; and \emph{prompt tuning}~\cite{li2021prefix,lester2021power}, which optimizes trainable prefix tokens attached to the model inputs. Among these options, model finetuning stands out for its conceptual simplicity and powerful performance, while importantly avoiding any additional inference-time delay.

The core idea behind model finetuning is to maintain a strong similarity between the pretrained model and its finetuned counterpart according to certain criteria, thereby safeguarding the pretrained knowledge. For example, many existing finetuning strategies employ a modest learning rate, which informally limits the divergence between the original and updated models. Considering the inherent structure within weight matrices, a more rigorous strategy aims to preserve their relational properties, specifically the angular relationships between neuron pairs. This motivates the Orthogonal Finetuning (OFT) framework~\cite{qiu2023controlling}, in which each layer’s neurons are collectively transformed via a single orthogonal matrix, ensuring that the angles between neuron vectors are exactly retained throughout finetuning. OFT has shown substantial potential for generalization and convergence in finetuning text-to-image diffusion models~\cite{qiu2023controlling}. Nonetheless, because orthogonal matrices are typically high-dimensional, OFT can incur a significant parameter overhead. To mitigate this, OFT incorporates a block-diagonal parameterization that limits the number of trainable parameters. Yet, this solution introduces new complications—namely, the fixed sparsity pattern inherent to block-diagonal matrices, with orthogonal operations isolated within each block. While this parameter reduction strategy increases efficiency, it also brings potential downsides: for instance, such block structures may restrict the model’s representational power and prevent them from approximating certain standard linear transformations. Critically, there remains an unresolved challenge in determining an optimal sparsity pattern for these orthogonal matrices.

The central challenge lies in constructing a dense orthogonal matrix without incurring excessive parameter costs. Ordinarily, forming a $d$-dimensional dense orthogonal matrix demands $\mathcal{O}(d^2)$ parameters, making a naive approach impractical. Our solution leverages a new method: synthesizing a dense orthogonal matrix as a product of several sparse, factorized matrices. This technique exploits the principle that storage requirements can be decreased at the expense of greater computational effort. Representing the dense orthogonal matrix as a series of sparse matrix multiplications necessitates performing these multiplications during every training iteration. To better contextualize matrix factorization, we recast the creation of a dense orthogonal matrix as analogous to a process of transmitting information through a network structured as a grid graph. Viewed through this lens, constructing a dense orthogonal matrix by successive sparse multiplications is tantamount to information propagation in such a graph. This transmission analogy enables us to devise numerous sparse matrix factorization schemes that control the parameter count while maintaining sufficient expressive capacity to yield dense orthogonal matrices.

In order to optimize parameter use within this information transmission setting, our approach is inspired by butterfly networks found in the Cooley-Tukey fast Fourier transform algorithm~\cite{cooley1965algorithm}, where efficient communication between distant nodes is possible~\cite{klasing1994broadcasting}. The butterfly graph architecture inherent in the Cooley-Tukey algorithm naturally suggests a type of sparse matrix factorization known as butterfly factorization. Employing butterfly factorization allows us to express a $d$-dimensional dense matrix as a product of $\mathcal{O}(\log d)$ sparse matrices, each containing only $\mathcal{O}(d)$ non-zero elements. By enforcing orthogonality on each sparse component, we obtain a highly efficient parameterization of orthogonal matrices, requiring only $\mathcal{O}(d \log d)$ parameters—a significant improvement over the full $\mathcal{O}(d^2)$ parameterization. Building upon this efficient structure, we introduce a new finetuning strategy, Orthogonal Butterfly~(BOFT), that unifies and generalizes earlier approaches such as OFT. Both block-diagonal and butterfly structures share the characteristic of inducing sparsity, thereby reducing the trainable parameter set within orthogonal matrices. In comparison, LoRA~\cite{hulora2022} employs low-rank matrices; thus, both low-rank and sparse matrices stand out as key classes of structured matrices~\cite{candes2011robust} for achieving parameter-efficient modeling.

In contrast to the block-diagonal form employed by OFT, which negotiates a balance between expressiveness and structural simplicity, BOFT leverages the butterfly architecture to achieve a more gradual transition between matrices from the complete orthogonal group (expressiveness) and the identity matrix (structural regularity). This approach narrows the space of hypotheses to a more compact subset within the orthogonal group, tailored for downstream tasks. Given that the butterfly structure underpins numerous rapid linear transformations—including the discrete Fourier and discrete cosine transforms—we posit that our structural design for parameter efficiency will inherently introduce inductive biases that enhance generalization and mitigate overfitting risks. This perspective is motivated by the butterfly structure’s capability to reproduce a range of standard linear transformations, a property unattainable for the block-diagonal form as used in OFT. Our main contributions are as follows:

\begin{itemize}[leftmargin=*,nosep]
    \item We tackle the issue of parameter-efficient orthogonal finetuning by introducing a novel information transmission framework, which conceptualizes the challenge of constructing a parameter-efficient dense orthogonal matrix as an information flow problem over a graph arranged in a grid topology.
    \item Drawing inspiration from butterfly networks found in the Cooley-Tukey algorithm, we introduce Orthogonal Butterfly, a new, parameter-efficient method for orthogonal finetuning within the context of the proposed information transmission framework.
    \item We offer theoretical justifications for the substantial parameter savings afforded by BOFT, alongside its strong expressive capacity and training stability. BOFT's matrix decomposition also enables unique forms of weight interpolation, as depicted in Figure~\ref{fig:boft_interpolation}.
    \item We present the first application of orthogonal finetuning across a diverse array of tasks beyond controllable text-to-image generation~\cite{qiu2023controlling}, highlighting BOFT's effectiveness as a general finetuning strategy. Specifically, we deploy BOFT in several adaptation scenarios spanning both computer vision and natural language processing, where it achieves marked improvements over current leading techniques, thereby demonstrating its superior parameter efficiency and generalization performance.
\end{itemize}

\section{Related Work}

\textbf{Parameter-efficient finetuning (PEFT)}. The surge in size and capabilities of foundation models (\emph{e.g.}, \cite{brown2020language,radford2021learning,rombach2022high,kirillov2023segment}) has heightened the demand for efficient methods to adapt such models for specific tasks without altering a large number of parameters~\cite{treviso2023efficient,ding2023parameter,lialin2023scaling}. Among the broad spectrum of PEFT strategies~\cite{houlsby2019parameter,aghajanyan2020intrinsic,hulora2022,edalati2022krona,wang2022adamix,gheini2021cross,zaken2022bitfit,guo2020parameter,sung2021training,ansell2022composable,lester2021power,li2021prefix,vu2022spot,he2021towards,mao2021unipelt,karimi2021compacter,liu2022few,sung2022lst,chen2022parameter,jia2022visual,chen2022adaptformer,zhang2022neural,jie2023fact,lian2022scaling,luo2023towards}, approaches that leverage reparameterization~\cite{aghajanyan2020intrinsic,hulora2022,edalati2022krona} align closely with our methodology. For example, LoRA~\cite{hulora2022} augments pretrained weight matrices by incorporating the product of two low-rank matrices, attaining strong results on tasks involving natural language. While LoRA constrains all inserted matrices to have a fixed rank, adaptive-rank variants~\cite{zhang2022adaptive,valipour2022dylora,zhang2023increlora} allow layer-specific ranks to better distribute the parameter constraints. The widespread adoption of this simple low-rank reparameterization underscores its practical effectiveness~\cite{dettmers2023qlora,zi2023delta,chavan2023one}. Drawing from the understanding that hyperspherical energy relates to generalization~\cite{liu2018learning,liu2021orthogonal}, \cite{qiu2023controlling} introduced orthogonal finetuning (OFT), which optimizes an orthogonal transformation matrix over the neurons within a layer, leading to better generalization and training stability relative to LoRA. Nevertheless, OFT often requires more trainable parameters compared to LoRA, creating motivation to enhance its parameter-efficiency. Additionally, the versatility of OFT for adaptation beyond text-to-image diffusion model control~\cite{qiu2023controlling} has not been established. By integrating butterfly factorization, BOFT addresses the parameter-efficiency of OFT. This advancement enables a more comprehensive demonstration of orthogonal finetuning's effectiveness across broader adaptation tasks.

\textbf{Butterfly structures}. The radix-2 Cooley-Tukey algorithm decomposes the $N$-point discrete Fourier transform into a recursive series of $\frac{N}{2}$-point transforms, creating a butterfly pattern expressible as a sequence of sparse matrix products, often called butterfly matrices. These matrices have a range of established applications: they help to represent orthogonal matrices in ways that eliminate the need for pivoting in Gaussian elimination~\cite{parker1995random}, boost the robustness of recurrent neural network training~\cite{jing2017tunable}, and serve in kernel approximation~\cite{munkhoeva2018quadrature}. Further, efficient linear operations have been developed by parameterizing transformations with butterfly matrices~\cite{dao2019learning,dao2019kaleidoscope}, and recent studies have leveraged butterfly matrices for scaling neural network training efficiency~\cite{chen2022pixelated,dao2022monarch}. Additionally, butterfly structures are instrumental in tasks such as rapid matrix-vector products~\cite{michielssen1996multilevel,o2010algorithm}, compressing matrix representations~\cite{li2015butterfly}, and enhancing network broadcasting~\cite{klasing1994broadcasting,li2003linear,rajasingh2016transmission}. Departing from these prior applications, we focus specifically on employing butterfly structures to improve the parameter efficiency of OFT.

\section{An Information Transmission View on Orthogonal Finetuning}

We begin by outlining foundational aspects of OFT. In order to adapt a pretrained weight matrix, OFT represents the new weight matrix as the product of a trainable orthogonal matrix and the fixed pretrained weight matrix. Unlike LoRA, which adjusts weights by adding a learnable low-rank matrix, OFT utilizes a multiplicative approach involving an orthogonal matrix. LoRA's parameter-efficiency is achieved through a low-rank parameterization, whereas the original OFT attains efficiency by adopting a block-diagonal orthogonal matrix structure~\cite{qiu2023controlling}; the smaller the diagonal blocks, the greater the efficiency in parameter usage. Figure~\ref{fig:comp_lora} provides a visual comparison between these methods. The rationale for employing orthogonal transformations during weight matrix finetuning is to maintain the pairwise angular relationships between neurons~\cite{liu2018learning,liu2021orthogonal,qiu2023controlling}, which helps to largely retain the semantic information from pretraining.

Specifically, for a pretrained linear layer $\bm{W}^0\in\mathbb{R}^{d\times n}$, OFT learns an orthogonal matrix $\bm{R}\in\mathbb{R}^{d\times d}$ and modifies the standard forward computation from $\bm{z}=(\bm{W}^0)^\top\bm{x}$ to $\bm{z}=(\bm{R}\bm{W}^0)^\top\bm{x}$, with $\bm{x}\in\mathbb{R}^{d}$ as the input vector and $\bm{z}\in\mathbb{R}^{n}$ as the output. To ensure the adaptation begins without altering pretrained knowledge, OFT initializes $\bm{R}$ as the identity matrix. Following the method described in~\cite{qiu2023controlling,liu2021orthogonal}, OFT maintains the orthogonality constraint on $\bm{R}$ during finetuning through the Cayley parameterization: $\bm{R}=(\bm{I}+\bm{Q})(\bm{I}-\bm{Q})^{-1}$, with $\bm{Q}$ being a skew-symmetric matrix such that $\bm{Q}=-\bm{Q}^\top$. For enhanced parameter-efficiency, $\bm{R}$ is constructed as a block-diagonal matrix $\text{diag}(\bm{R}_1,\bm{R}_2,\cdots,\bm{R}_r)$, where each block $\bm{R}_i\in\mathbb{R}^{b\times b}$ is an orthogonal matrix and $br=d$ holds. This block-diagonal approach reduces parameter count but introduces a structural assumption: it implicitly partitions the dimensions of each neuron (i.e., columns of $\bm{W}^0$) into $r$ distinct groups, with each group being independently transformed via separate orthogonal matrices. While empirical results support the efficacy of this approach, the grouping of neuron dimensions by simple index partition lacks principled justification, thereby motivating interest in fully dense orthogonal matrices. Consequently, a key question arises: \emph{Is it possible to design a dense orthogonal matrix formulation that retains the parameter-efficiency of block-diagonal structures?}

To tackle this issue, we suggest representing a dense orthogonal matrix as the product of several sparse orthogonal matrices. To systematically investigate the sparsity structures arising from orthogonal matrix factorization, we reinterpret the task of constructing a dense orthogonal matrix in OFT within the framework of an information transmission process. Concretely, assembling a dense matrix $\bm{R}\in\mathbb{R}^{d\times d}$ via the product $\bm{R}=\bm{B}_m\bm{B}_{m-1}\cdots\bm{B}_1$, where each $\bm{B}_i$ is a $d \times d$ matrix, can be seen as a process of information propagation through a $d\times (m+1)$ network grid, as depicted in Figure~\ref{fig:info_trans}. This analogy is motivated by the perspective that a dense $d$-dimensional square matrix essentially encodes all-to-all connections from one set of $d$ nodes to another. For $\bm{R}$, an entry $R_{ij}=0$ signifies the absence of a communication path from node $j$ to node $i$, whereas a nonzero $R_{ij}$ allows for transmission. Therefore, expressing $\bm{R}$ as a product $\bm{B}_m\bm{B}_{m-1}\cdots\bm{B}_1$ is analogous to a sequence of information exchanges governed by the connectivity patterns of the intermediate matrices $\bm{B}_i$. Here, the flow of information starts according to $\bm{B}_1$ and is concluded by $\bm{B}_n$. 

An explicit instance is provided in Figure~\ref{fig:info_trans}, which illustrates the factorization $\bm{R}=\bm{B}_5\bm{B}_4\bm{B}_3\bm{B}_2\bm{B}_1$, along with its respective sparsity patterns and induced graph. The visualized graph represents the full unrolling of matrix multiplication across layers. Each matrix $\bm{B}_i$ defines how nodes at layer $i$ are connected to those at layer $i+1$: the entry $(j_1, j_2)$ in $\bm{B}_i$ indicates the presence (or absence, if zero) of a directed edge from the $j_2$-th node at layer $i$ to the $j_1$-th node at layer $i+1$. For the overall product $\bm{B}_5\bm{B}_4\bm{B}_3\bm{B}_2\bm{B}_1$ to yield a dense $\bm{R}$, it is required that each initial (level-one) node has a path to every node at the sixth level. As a specific case, considering just $\bm{R}=\bm{B}_4\bm{B}_3\bm{B}_2\bm{B}_1$—which restricts our focus to transmission from first-level sources to fifth-level receivers—shows that the first node at the source cannot communicate with the third node at the receiver side; this manifests as a zero at position $(3,1)$ in the product's sparsity pattern. Analogously, the interplay between induced graphs and sparsity holds when considering shorter products like $\bm{R}=\bm{B}_3\bm{B}_2\bm{B}_1$ and $\bm{R}=\bm{B}_2\bm{B}_1$. In general terms, ensuring that $\bm{R}\in\mathbb{R}^{d\times d}$ is dense means the factorization matrices $\bm{B}_m,\ldots,\bm{B}_1$ must specify a collection of directed edges on a $d\times (m+1)$ grid. Each directed edge links two nodes from consecutive levels (that is, adjacent columns), guaranteeing that every initial node can reach all nodes in the final layer through some sequence of connections.

Adopting the information transmission perspective enables a deeper analysis of the sparsity structures in the factorization matrices used in OFT. As depicted in Figure~\ref{fig:blk_diag}, the original OFT generates a block-diagonal form for $\bm{R}$. Although this structure diminishes the count of trainable parameters, it inherently restricts the construction of a fully dense matrix $\bm{R}$. Our objective is thus to assemble a dense orthogonal matrix by leveraging $m$ sparse orthogonal matrices while employing a minimal set of essential trainable parameters. From the standpoint of information transmission, two primary requirements emerge: \emph{(i)} \textbf{dense connectivity}, wherein each node at the initial level maintains at least one route to every terminal node, and \emph{(ii)} \textbf{minimum free edges}, ensuring that the aggregate number of connections is minimized subject to the orthogonality constraint. It is important to emphasize that orthogonality enforces subtle restrictions on edge formation between adjacent layers. For instance, to ensure each matrix $\bm{B}_i$ retains full rank—an essential property for orthogonality—it is required to establish $d$ connections that create a bijective mapping between all nodes at level $i$ and those at level $(i=1)$. This sets a lower bound of $d$ connections between successive layers (\emph{e.g.}, in the example in Figure~\ref{fig:info_trans}, $4$ edges are needed). These $d$ mandatory connections uphold orthogonality but are non-trainable and thus excluded from the edge count, as their entries are restricted to fixed values such as $\pm 1$ in a $d\times d$ orthogonal matrix. The orthogonality property further reduces the parameter count compared to general dense matrices, introducing extra dependencies between edges. For instance, in a $2\times2$ orthogonal matrix, assigning a zero to the $(1,1)$ position necessitates a zero at the $(2,2)$ position as well (\emph{i.e.}, omitting one connection can enforce the omission of another). Consequently, for any chosen edge configuration, the orthogonality requirement may necessitate adjustment—sometimes adding or omitting further edges. The information transmission framework, focusing on the patterns of nonzero entries in the resultant orthogonal matrix, offers considerable value for OFT, as our interest centers solely on which elements of $\bm{R}$ are trainable, independent of their numerical values.

A straightforward fully connected scheme between two layers requires $\mathcal{O}(d^2)$ connections (\emph{i.e.}, one dense orthogonal matrix), resulting in $d^2 - d$ total connections (for example, with $d=4$, there are 12 connections). Figure~\ref{fig:info_trans} illustrates a matrix factorization that is achievable, which utilizes only $10$ connections—fewer than a dense orthogonal alternative. This graphical approach facilitates the analysis of OFT’s parameter efficiency, as it allows for easy construction of various feasible factorizations. Drawing on the network topology used in the Cooley-Tukey algorithm, namely the butterfly graph~\cite{cooley1965algorithm}, we note that this structure enables efficient all-to-all connectivity between $d$ input and $d$ output nodes, requiring only $\mathcal{O}(d\log d)$ connections. As a concrete example, the configuration in Figure~\ref{fig:info_trans} connects $d$ nodes with $10$ edges, yet a butterfly topology achieves dense connections with just $8$ edges. The following discussion explains how adopting the butterfly network further enhances parameter efficiency.

\section{Orthogonal Parameterization by Butterfly Factorization}

Originally, the butterfly structure was introduced in the Cooley-Tukey algorithm for fast Fourier transform computations. In this context, a local adjustment in frequency space can lead to global effects in the spatial domain—a property akin to the information flow challenge under consideration, where each node at the initial layer is able to communicate with all nodes at the terminal layer. Beyond Fourier transform, butterfly networks are also widely used in computer systems as a topology that supports effective and scalable data exchange~\cite{solihin2015fundamentals,li2011linear}. When $k\geq 2$ and is a power of two, we define the \emph{butterfly factor} $\bm{B}^F(k)$ as follows:

\begin{equation}\label{eq:but_factor}
\begin{aligned}
    \bm{B}^F(k)=\begin{bmatrix}
    \text{diag}(\bm{d}_1) & \text{diag}(\bm{d}_2) \\
    \text{diag}(\bm{d}_3) & \text{diag}(\bm{d}_4)
    \end{bmatrix}\in\mathbb{R}^{k\times k},
\end{aligned}
\end{equation}

where each $\bm{d}_i\in\mathbb{R}^{\frac{k}{2}}$ for $i=1,2,3,4$ are vectors. With $d=2^N$, we recursively define the $d$-dimensional \emph{butterfly matrix} $\bm{B}(d)\in\mathbb{R}^{d\times d}$ by

\begin{equation}
\begin{aligned}
    \!\!\bm{B}(d)=\tilde{\bm{B}}(d,d)\!\cdot\!\begin{bmatrix}
    \bm{B}_1(\frac{d}{2})\!\!\!\!\! & \bm{0} \\
    \bm{0}\!\!\!\!\! &\bm{B}_2(\frac{d}{2})
    \end{bmatrix}= \tilde{\bm{B}}(d,d)\tilde{\bm{B}}(d,\frac{d}{2})\cdots\tilde{\bm{B}}(d,2),
\end{aligned}
\end{equation}

Here, $\bm{B}_1(\frac{d}{2})$ and $\bm{B}_2(\frac{d}{2})$ denote two butterfly matrices each of dimension $\frac{d}{2}$. We introduce the \emph{butterfly component} as $\thickmuskip=2mu \medmuskip=2mu \tilde{\bm{B}}(d,k)=\text{diag}(\bm{B}^F_1(k),\cdots,\bm{B}^F_{d/k}(k))$, which forms a block-diagonal matrix of dimension $d\times d$ with blocks of size $k$. Here, each $\bm{B}^F_i(k)$, for every $i$, corresponds to a butterfly factor as detailed in Equation~\ref{eq:but_factor}. With these preliminaries, the butterfly matrix can be employed for orthogonal matrix parameterization. To ensure the resulting matrix is orthogonal, it suffices that every multiplicative term $\tilde{\bm{B}}(d,k)$ in the butterfly decomposition $\bm{B}(d)$ is orthogonal. Focusing first on $\tilde{\bm{B}}(d,2)$, a block-diagonal matrix with $2\times 2$ blocks, we can readily ensure orthogonality by using the Cayley transform (or equivalently, $2$-dimensional rotation) for parameterizing each block, in line with the approaches from \cite{liu2021orthogonal,qiu2023controlling}. Each butterfly component essentially represents a permutation of the sparsity pattern found in $\tilde{\bm{B}}(d,2)$; thus, all butterfly components can be straightforwardly parameterized to yield orthogonal matrices. This results in a scalable approach for constructing orthogonal matrices using multiple $2\times 2$ orthogonal blocks. Extending this approach as in \cite{chen2022pixelated}, we define the block butterfly component $\tilde{\bm{B}}^b(d,k)$, where every entry in $\bm{d}_i$ for all $i$ is promoted to a $b\times b$ matrix. Orthogonality of the block butterfly component $\tilde{\bm{B}}^b(d,2)$ can be guaranteed by parameterizing each $2b\times 2b$ block to be orthogonal. For other butterfly components $\tilde{\bm{B}}^b(d,k)$ with $k>2$, the sparsity structure is a block-wise permutation of that in $\tilde{\bm{B}}^b(d,2)$, and hence, these can also be parameterized as orthogonal matrices using the same strategy. Altogether, the forward propagation in BOFT is described as

\begin{equation}
    \begin{aligned}\nonumber
    \bm{z}=\big(\bm{R}(m,b)\cdot\bm{W}^0\big)^\top\bm{x},~~~~\text{s.t.}~\bigg{\{}\bm{R}(m,b)=\prod_{i=1}^m \tilde{\bm{B}}^b_{(d,i)}~~\&~~\underbrace{\big(\tilde{\bm{B}}^b_{(d,j)}\big)^\top\tilde{\bm{B}}^b_{(d,j)}=\tilde{\bm{B}}^b_{(d,j)}\big(\tilde{\bm{B}}^b_{(d,j)}\big)^\top\!=\bm{I}_d}_{\forall j\in [1, m]}\bigg{\}},
    \end{aligned}
\end{equation}

Here, we use the notation $\tilde{\bm{B}}^b(d,2^{m - i+1})$ as shorthand for $\tilde{\bm{B}}^b_{(d,i)}$, and let $\bm{I}_d$ represent the $d$-dimensional identity matrix. The orthogonal transformation $\bm{R}(m,b)\in\mathbb{R}^{d\times d}$ is constructed by multiplying $m$ orthogonal butterfly matrices. To simplify notation, we refer to BOFT that uses $\bm{R}(m,\frac{b}{2})$ as BOFT($m$,$b$), under the constraint $b\geq 2$. When $m=1$, BOFT($1$,$b$) simplifies to the block-diagonal OFT~\cite{qiu2023controlling} with block dimension $b$. Similarly, if BOFT($1$,$d$) is employed, the setup recovers the original OFT~\cite{qiu2023controlling} utilizing an unconstrained, full-size orthogonal matrix. Choosing $m=\log \frac{2d}{b}$ and $b$ as the block size results in BOFT($\log \frac{2d}{b}$,$b$), where the orthogonal matrix $\bm{R}$ is constructed as a dense butterfly matrix $\bm{B}^{\frac{b}{2}}(d)$. More broadly, BOFT($m$,$b$) contains $\frac{1}{2}(b-1)dm$ tunable parameters when adapting a linear module of shape $d\times n$. Setting $m=\log d$ and $b=2$, so that BOFT becomes a classic butterfly matrix, the parameter complexity becomes $\mathcal{O}(d\log d)$. In comparison, the standard OFT—using a dense orthogonal parameterization—requires $\mathcal{O}(d^2)$ parameters, and a block-diagonal OFT with $r$ blocks incurs $\mathcal{O}(bd)$. Thus, for dense orthogonality, traditional OFT must select $b=d$ to fully utilize its parameterization, while BOFT permits arbitrary $b$ to attain the same outcome.

\textbf{Identity initialization for BOFT}. Conventional fine-tuning approaches are initialized to mirror the pretrained weights to ensure minimal deviation at the beginning of adaptation; for example, LoRA initializes its low-rank matrices with zeros. Analogously, BOFT initializes all butterfly matrices as identity matrices, which is realized by setting their skew-symmetric generator matrices to zeros under the Cayley transform.

\textbf{Multiplicative Dropout for BOFT}. LoRA~\cite{hulora2022} leverages a Dropout mechanism for its low-rank updates to regularize against overfitting. While standard Dropout~\cite{srivastava2014dropout} applies straightforwardly to LoRA’s additive formulation, it is incompatible with BOFT’s multiplicative structure. To this end, we introduce a multiplicative Dropout tailored to BOFT: since $\bm{R}(m,b)$ comprises $m$ butterfly blocks, each reducible to $2b\times 2b$ block-diagonal orthogonal matrices, our Dropout randomly selects a proportion $p_1$ of these butterfly blocks and, within each, a proportion $p_2$ of the diagonal blocks, substituting the selected blocks with identity matrices.

\section{Intriguing Insights and Discussions}

\textbf{Expressivity of BOFT}. While the butterfly structure, combined with permutations, has been shown to exactly replicate numerous well-known fast linear transforms~\cite{dao2019learning,dao2019kaleidoscope} (\emph{e.g.}, fast Fourier transform, Hadamard transform), its ability to approximate arbitrary orthogonal matrices is less well-understood. To investigate this, we simulate the approximation of a randomly generated dense orthogonal matrix~\cite{anderson1987generation} of size $512\times 512$. The outcomes, displayed in Figure~\ref{fig:para_eff}, are averaged across 10 random initializations. Here, the y-axis captures the approximation error, while the x-axis indicates the count of trainable parameters. Each color corresponds to BOFT models using identical block sizes; the point at the far left represents the error of BOFT($1$,$b$), which is equivalent to the basic block-diagonal OFT with block size $b$. Overall, BOFT demonstrates superior parameter efficiency compared to OFT. Notably, BOFT($9$,$2$) achieves greater expressivity than BOFT($1$,$16$) while employing significantly fewer parameters. Employing smaller $b$ alongside larger $m$ in BOFT configurations generally improves parameter efficiency; for example, BOFT($6$,$4$) maintains a similar approximation error as BOFT($2$,$16$) yet uses substantially fewer parameters. These findings support the notion that the butterfly matrix structure captures a richer subset of the orthogonal group relative to block-diagonal counterparts, thus making BOFT provably more expressive than OFT at equal block sizes.

\begin{theorem}[Expressivity of BOFT]\label{thm:exp_boft}
    BOFT offers greater expressive power than OFT for a given block size. Specifically, to approximate any orthogonal matrix of size $d$, one may construct products of butterfly matrices as follows: $\bm{B}_{d-1,1}(d)\bm{B}^\top_{d-1,2}(d)\cdots\bm{B}_{1,1}(d)\bm{B}^\top_{1,2}(d)$, where each $\bm{B}_{i,j}(d),\forall i,\forall j$ is a butterfly matrix.
\end{theorem}

Theorem~\ref{thm:exp_boft} points toward a natural extension for BOFT—namely, generalizing the final orthogonal matrix via the composition $\thickmuskip=2mu \medmuskip=2mu \bm{R}^G(m_1,b_1,m_2,b_2,l)=\bm{R}_{l,1}(m_1,b_1)\bm{R}_{l,2}^T(m_2,b_2)\cdots\bm{R}_{1,1}(m_1,b_1)\bm{R}_{1,2}^T(m_2,b_2)$, where $\bm{R}_{i,j}^T(m,b)$ is an orthogonal matrix of the BOFT framework. In the case where $m_1=m_2=\log d$, $b_1=b_2=2$, and $l=d-1$, $\bm{R}^G(m_1,b_1,m_2,b_2,l)$ is capable of spanning the full orthogonal group—a composition also known as the kaleidoscope hierarchy~\cite{dao2019kaleidoscope}. Nevertheless, it is important to acknowledge that higher expressivity does not inherently translate to better finetuning performance; even with universal expressiveness, full finetuning can sometimes result in disappointing outcomes. Ultimately, the balance between expressiveness and structural regularity governs model generalization in finetuning tasks. By incorporating structural priors, BOFT expands the available finetuning parameter space, facilitating a more effective compromise between these two facets.

\textbf{Spectral properties}. Orthogonal finetuning tends to outperform LoRA regarding spectral characteristics because it maintains the spectral norm of the original pretrained weight matrix $\bm{W}^0$ without alteration. To illustrate, applying the singular value decomposition yields $\bm{W}^0 = \bm{U}\bm{\Sigma}\bm{V}^\top$, where $\bm{U}$ and $\bm{V}$ are orthogonal matrices, and $\bm{\Sigma}$ contains the singular values along its diagonal. In both OFT and BOFT, the finetuning process introduces a left-multiplied orthogonal matrix $\bm{R}$, producing $\bm{R}\bm{U}\bm{\Sigma}\bm{V}^\top$ as the new weight matrix; importantly, this operation leaves the largest singular value (\emph{i.e.}, the spectral norm of $\bm{W}^0$) unchanged. Such invariance in the spectral norm has been empirically demonstrated to promote both stability during optimization and enhanced generalization~\cite{miyato2018spectral, yoshida2017spectral}. Additional mathematical analysis of these properties is provided in Appendix~\ref{app:sn_property}.

\textbf{Orthogonal finetuning as learning bilinear similarity}. BOFT admits a formulation as learning the bilinear similarity $\bm{w}^0_i\bm{R}\bm{x}$, in which $\bm{w}^0_i$ corresponds to the $i$-th column (neuron) of $\bm{W}^0$. From this perspective, BOFT is equivalent to estimating the bilinear similarity matrix $\bm{R}$, subject to the orthogonality constraint—a strong form of regularization. This connection naturally relates BOFT to the frameworks of distance metric learning~\cite{xing2002distance} and bilinear forms~\cite{roman2005advanced}.

\textbf{Inductive bias and generalization in BOFT}. In BOFT, $\bm{R}(m,b)$ typically embodies a subset of the orthogonal group characterized by specific structural constraints, thereby narrowing the range of possible solutions and imparting an inductive bias. We posit that the inductive bias imposed through butterfly factorization aids in generalization, since this structure recurs in a range of classical linear operators~\cite{dao2019kaleidoscope}, including the discrete Fourier, sine/cosine, and Hadamard transforms. Additionally, the use of sparse factorization within BOFT may confer further implicit inductive biases~\cite{gunasekar2017implicit, li2018algorithmic, liu2019neural}.

\textbf{Relation to butterfly-based sparse training}. Several studies~\cite{dao2019learning,dao2019kaleidoscope,chen2022pixelated,dao2022monarch} have explored sparse training paradigms grounded in butterfly parameterizations. Typically, these approaches directly reparameterize neural network weight matrices using butterfly structures and perform training from initialization. The work in~\cite{dao2022monarch} extends this line by finetuning pretrained networks: weights are first projected onto a modified butterfly matrix space and subsequent optimization occurs on these projected elements for specific downstream tasks. In contrast, BOFT introduces a distinctive finetuning mechanism, utilizing layer-consistent transformation matrices applied to the model weights.

\input{rebuttal-results}

\section{Concluding Remarks and Limitations}

This work introduces Orthogonal Butterfly, a broadly applicable method for parameter-efficient finetuning of foundation models utilizing the butterfly pattern. Our central contribution lies in achieving improved parameter efficiency by expressing dense orthogonal matrices as products of several sparse orthogonal matrices. To facilitate practical identification of such matrix decompositions, we present a graphical information transmission formalism. Employing this formalism reveals that butterfly matrices can reliably provide sparse orthogonal matrix factorization as desired. Empirical evaluations substantiate BOFT’s effectiveness for finetuning across diverse model families, including large language models, vision architectures, and text-to-image generators. These experiments confirm that BOFT delivers superior performance as a general-purpose finetuning approach.

Notwithstanding its empirical advantages, BOFT has certain shortcomings. Specifically, due to representing the final orthogonal transformation as a product of multiple matrices, its training introduces marginally higher computational overhead than OFT. Addressing the training efficiency of BOFT thus represents an open avenue for future work. Encouragingly, once finetuning is complete, the orthogonal matrices derived by BOFT can be merged directly into the pretrained model, ensuring that inference speed is not impacted. In addition, it remains unclear whether the butterfly topology provides optimal efficiency for information transmission. The information transmission framework established in this work also facilitates cross-pollination with insights from computer networking, a field extensively examining topologies for efficient information delivery. We anticipate that alternative network architectures, inspired by such studies, could yield more effective compositions of dense orthogonal matrices.

\end{document}